% !PRINTED 499
\weberpage{499}
Denn bezeichnen wir, wenn $\alpha$ irgend eine der Ziffern
$0, 1, 2, \ldots n - 1$ bedeutet, mit $\alpha_\pi$ und $\alpha_{\pi'}$ die Ziffern, in die $\alpha$
durch $\pi$ und durch $\pi'$ übergeht, so ist $\alpha_{\pi\pi'}$ die Ziffer, in die $\alpha_\pi$
durch $\pi'$ oder $\alpha$ durch $\pi\pi'$ übergeht. Sind $\alpha, \beta, \gamma \ldots$ die Ziffern
aus dem System $M$, so ist
\[
\begin{gathered}
\pi'^{-1}(\alpha, \beta) \pi' = (\alpha_{\pi'}, \beta_{\pi'}) \\
\pi\pi'^{-1}(\alpha, \beta) \pi'\pi^{-1} = (\alpha_{\pi'\pi^{-1}}, \beta_{\pi'\pi^{-1}}) = (\gamma, \delta).
\end{gathered}
\]
Ist nun $\alpha_{\pi'} = \gamma_\pi$ in $M'$ enthalten, so ist $\alpha_{\pi'\pi^{-1}} = \gamma$ in $M$
enthalten, und folglich muss $\beta_{\pi'\pi^{-1}} = \delta$ gleichfalls in $M$ enthalten sein,
da die Transposition $(\gamma, \delta)$ in $P$ vorkommt. Folglich ist $\alpha_{\pi'} = \gamma_\pi$
und $\beta_{\pi'} = \delta_\pi$, d.\,h. wenn eine der Ziffern aus $M''$ in $M'$ enthalten ist, so
müssen $M'$ und $M''$ identisch sein.
Es zerfällt also die Gesammtheit der Ziffern $0, 1 \ldots n - 1$
in Systeme $M, M', M'' \ldots$ von je $m$ Ziffern, die durch $P$ imprimitiv permutirt
werden, d.\,h. die Gruppe $P$ ist imprimitiv.
Wenn $n$ eine Primzahl ist, so giebt es keine imprimitiven Permutationsgruppen;
also ist in diesem Falle eine transitive Gruppe, die eine Transposition enthält,
nur die symmetrische Gruppe.
Ganz auf dieselbe Weise schliessen wir aus dem Satze 7:

10. \gesperrt{Wenn eine transitive Permutationsgruppe von
$n$ Ziffern einen dreigliedrigen Cyklus enthält,
so enthält sie die ganze alternirende Gruppe
oder sie ist imprimitiv.}

Denn enthält die Gruppe $P$ die dreigliedrigen Cyklen
\[
(0, 1, 2), (0, 1, 3) \ldots (0, 1, m - 1),
\]
aber keinen anderen mit den Ziffern $0, 1$, so enthält sie nach 7.
die ganze alternirende Gruppe der Ziffern $M = 0, 1, 2 \ldots m - 1$.
Ist nun $m < n$, so kann in $P$ kein dreigliedriger Cyklus vorkommen, der eine
oder zwei der Ziffern von $M$ mit anderen Ziffern verbindet. Denn enthält $P$
z.\,B. $(0, m, m + 1)$, so enthält es, wie man aus den Zusammensetzungen
\[
(0, 2, 1) (0, m, m + 1) (0, 1, 2) = (1, m, m + 1)
\]
ersieht, auch $(1, m, m + 1), \ldots (m - 1, m, m + 1)$ und also
nach 7. die ganze alternirende Gruppe der Ziffern $0, 1, 2 \ldots m + 1$,
was gegen die Voraussetzung ist. Es zerfallen also die $n$ Ziffern
von $P$, wie vorhin, in mehrere Systeme $M, M', M'' \ldots$ der Imprimitivität. Wir können noch hinzufügen,
%CONT>>
